
\documentclass[12pt, a4paper]{article}

% ─── Codificación y fuente ───────────────────────────────────────────────────
\usepackage[utf8]{inputenc}
\usepackage[T1]{fontenc}
\usepackage[spanish]{babel}
\usepackage{lmodern}

% ─── Matemáticas ────────────────────────────────────────────────────────────
\usepackage{amsmath, amssymb, amsthm}

% ─── Geometría de página ─────────────────────────────────────────────────────
\usepackage[top=2.5cm, bottom=2.5cm, left=3cm, right=3cm]{geometry}

% ─── Gráficos y colores ──────────────────────────────────────────────────────
\usepackage{graphicx}
\usepackage[dvipsnames, table]{xcolor}
\definecolor{azuloscuro}{RGB}{0, 51, 102}
\definecolor{azulclaro}{RGB}{220, 235, 252}
\definecolor{grisclaro}{RGB}{245, 245, 245}
\definecolor{verdeos}{RGB}{0, 100, 60}
\definecolor{rojosuave}{RGB}{160, 30, 30}

% ─── Hipervínculos ───────────────────────────────────────────────────────────
\usepackage[hidelinks, colorlinks=true, linkcolor=azuloscuro,
            urlcolor=azuloscuro, citecolor=azuloscuro]{hyperref}

% ─── Cajas y entornos destacados ─────────────────────────────────────────────
\usepackage{tcolorbox}
\tcbuselibrary{skins, breakable, theorems}

\tcbset{
  hipotesis/.style={
    colback=azulclaro, colframe=azuloscuro,
    fonttitle=\bfseries\small, title={Hipótesis},
    breakable, before skip=8pt, after skip=8pt,
    left=6pt, right=6pt
  },
  nota/.style={
    colback=grisclaro, colframe=gray!60,
    fonttitle=\bfseries\small, title={Nota},
    breakable, before skip=8pt, after skip=8pt,
    left=6pt, right=6pt
  },
  clave/.style={
    colback=yellow!10, colframe=orange!70!black,
    fonttitle=\bfseries\small, title={Idea clave},
    breakable, before skip=8pt, after skip=8pt,
    left=6pt, right=6pt
  },
  teorema/.style={
    colback=green!5, colframe=verdeos,
    fonttitle=\bfseries\small,
    breakable, before skip=8pt, after skip=8pt,
    left=6pt, right=6pt
  }
}

% ─── Código fuente ───────────────────────────────────────────────────────────
\usepackage{listings}
\lstset{
  backgroundcolor=\color{grisclaro},
  basicstyle=\ttfamily\small,
  keywordstyle=\color{azuloscuro}\bfseries,
  commentstyle=\color{gray},
  stringstyle=\color{rojosuave},
  breaklines=true,
  frame=single,
  rulecolor=\color{gray!50},
  numbers=left,
  numberstyle=\tiny\color{gray},
  numbersep=6pt,
  tabsize=4,
  showstringspaces=false,
  captionpos=b,
  literate=
    {á}{{\'a}}1 {é}{{\'e}}1 {í}{{\'{\i}}}1
    {ó}{{\'o}}1 {ú}{{\'u}}1 {ñ}{{\~n}}1
    {Á}{{\'A}}1 {É}{{\'E}}1 {Í}{{\'I}}1
    {Ó}{{\'O}}1 {Ú}{{\'U}}1 {Ñ}{{\~N}}1
    {←}{{$\leftarrow$}}1 {≠}{{$\neq$}}1
}

% ─── Tablas ──────────────────────────────────────────────────────────────────
\usepackage{booktabs}
\usepackage{array}
\usepackage{longtable}
\newcolumntype{L}[1]{>{\raggedright\arraybackslash}p{#1}}
\newcolumntype{C}[1]{>{\centering\arraybackslash}p{#1}}

% ─── Encabezados ─────────────────────────────────────────────────────────────
\usepackage{fancyhdr}
\pagestyle{fancy}
\fancyhf{}
\fancyhead[L]{\small\textcolor{azuloscuro}{\textbf{Monte Carlo Tree Search}}}
\fancyhead[R]{\small\textcolor{gray}{Introducción desde cero}}
\fancyfoot[C]{\thepage}
\renewcommand{\headrulewidth}{0.4pt}

% ─── Tipografía de títulos ───────────────────────────────────────────────────
\usepackage{titlesec}
\titleformat{\section}
  {\large\bfseries\color{azuloscuro}}{\thesection.}{0.7em}{}[\vspace{-2pt}\rule{\linewidth}{0.4pt}]
\titleformat{\subsection}
  {\normalsize\bfseries\color{azuloscuro!80}}{\thesubsection.}{0.5em}{}
\titleformat{\subsubsection}
  {\normalsize\bfseries\color{gray!80!black}}{\thesubsubsection.}{0.5em}{}

% ─── Párrafo ─────────────────────────────────────────────────────────────────
\setlength{\parskip}{6pt}
\setlength{\parindent}{0pt}

% ─── Contador de hipótesis ────────────────────────────────────────────────────
\newcounter{hipnum}
\newcommand{\hip}[1]{%
  \stepcounter{hipnum}%
  \begin{tcolorbox}[hipotesis, title={Hipótesis H\thehipnum}]
    #1
  \end{tcolorbox}%
}

% ─────────────────────────────────────────────────────────────────────────────
\begin{document}

% ══════════════════════════════════════════════════════════════════════════════
%  PORTADA
% ══════════════════════════════════════════════════════════════════════════════
\begin{titlepage}
  \centering
  \vspace*{2cm}
  {\Huge\bfseries\color{azuloscuro} Monte Carlo Tree Search\par}
  \vspace{0.5cm}
  {\LARGE\color{gray!80!black} Introducción rigurosa desde cero\par}
  \vspace{1.5cm}
  \rule{0.6\linewidth}{1pt}
  \vspace{1cm}

  \begin{tcolorbox}[colback=azulclaro, colframe=azuloscuro, width=0.85\linewidth]
    \centering\small
    Este documento presenta el algoritmo Monte Carlo Tree Search (MCTS)
    partiendo de cero, enunciando explícitamente las hipótesis en cada paso,
    e incluyendo referencias a recursos externos curados por nivel.
  \end{tcolorbox}

  \vspace{1.5cm}
  {\large Contenidos:\par}
  \vspace{0.4cm}
  \begin{tabular}{cl}
    \textbullet & El problema de decisión secuencial \\
    \textbullet & El bandido multi-brazo y UCB1 \\
    \textbullet & Las cuatro fases de MCTS \\
    \textbullet & UCT: Upper Confidence Bounds for Trees \\
    \textbullet & Garantías teóricas de convergencia \\
    \textbullet & Pseudocódigo comentado \\
    \textbullet & Cuándo se cumplen/violan las hipótesis \\
    \textbullet & Recursos por nivel de profundidad \\
  \end{tabular}

  \vfill
  {\small\color{gray} Compilado con \LaTeX{} \textbullet{} Marzo 2026}
\end{titlepage}

\tableofcontents
\clearpage

% ══════════════════════════════════════════════════════════════════════════════
\section{¿Qué problema queremos resolver?}
% ══════════════════════════════════════════════════════════════════════════════

Imagina que juegas al ajedrez. En cada turno tienes que elegir un movimiento.
El problema es que el árbol de posibilidades futuras es astronómico: el ajedrez
tiene aproximadamente $10^{120}$ partidas posibles (el llamado \textbf{número
de Shannon}). Ningún computador puede explorar todo ese árbol en tiempo útil.

Durante décadas la solución fue \textbf{minimax con poda alfa-beta}: explorar
el árbol hasta cierta profundidad y usar una función de evaluación escrita a
mano (p.\,ej., «el caballo vale 3 peones»). Esto funcionó razonablemente bien
para ajedrez, pero fracasó en \textbf{Go}: el tablero es más grande, el factor
de ramificación es $\approx 250$, y nadie supo escribir una buena función de
evaluación para Go durante décadas.

MCTS nació precisamente para resolver este vacío:

\begin{tcolorbox}[clave]
¿Cómo tomar buenas decisiones cuando el árbol es demasiado grande para
explorarlo exhaustivamente \emph{y} no disponemos de una función de evaluación
de posiciones?
\end{tcolorbox}

La respuesta de MCTS es \textbf{muestrear el futuro aleatoriamente} (método
Monte Carlo) y \textbf{guiar la búsqueda} hacia las ramas más prometedoras
usando estadísticos acumulados. El resultado es un algoritmo que:

\begin{itemize}
  \item No necesita función de evaluación escrita a mano.
  \item Mejora su calidad cuanto más tiempo de cómputo se le da (\emph{anytime}).
  \item Converge a la decisión óptima en el límite.
\end{itemize}

% ══════════════════════════════════════════════════════════════════════════════
\section{Prerrequisito: el problema del bandido multi-brazo}
% ══════════════════════════════════════════════════════════════════════════════

Antes de entender MCTS es imprescindible entender el problema matemático que
reside en su corazón.

\subsection{Planteamiento}

\textbf{Situación:} Estás en un casino con $K$ máquinas tragaperras (los
«brazos» del bandido). Cada máquina $i \in \{1,\ldots,K\}$ tiene una
distribución de recompensa $\nu_i$ con media desconocida $\mu_i$. Puedes tirar
una palanca a la vez. Quieres \textbf{maximizar la recompensa total} en $n$
tiradas.

\hip{Las recompensas de cada máquina son \textbf{independientes entre sí e
idénticamente distribuidas} (i.i.d.) a lo largo del tiempo \emph{para esa
máquina}. Es decir, la distribución de la máquina $i$ no cambia con el tiempo
ni depende de las otras máquinas.}

\hip{Las recompensas están \textbf{acotadas} en el intervalo $[0,1]$ (o en
algún intervalo conocido $[a,b]$). Esta hipótesis es necesaria para obtener
cotas de concentración (desigualdad de Hoeffding) que garantizan los resultados
teóricos.}

El dilema fundamental se llama \textbf{equilibrio exploración--explotación}:

\begin{itemize}
  \item \textbf{Explotación:} tiramos siempre de la máquina que mejor nos ha
    ido hasta ahora.
  \item \textbf{Exploración:} probamos máquinas poco visitadas por si son
    mejores.
\end{itemize}

Si solo explotamos, podemos quedarnos atascados en un óptimo local. Si solo
exploramos, desperdiciamos tiradas en máquinas malas.

\subsection{El arrepentimiento (regret)}

Para medir la calidad de un algoritmo de bandido, definimos el
\textbf{arrepentimiento acumulado} tras $n$ tiradas:

\begin{equation}
  R_n \;=\; n\,\mu^* - \sum_{t=1}^{n} r_t
\end{equation}

donde $\mu^* = \max_i \mu_i$ es la recompensa de la mejor máquina y $r_t$ es
la recompensa obtenida en la tirada $t$. Un algoritmo es bueno si $R_n$ crece
lentamente con $n$.

\begin{tcolorbox}[teorema, title={Resultado fundamental (Lai \& Robbins, 1985)}]
Bajo las hipótesis H1 y H2, ningún algoritmo puede lograr regret menor que
$\Omega(\log n)$ en general. El regret logarítmico es por tanto \textbf{óptimo}.
\end{tcolorbox}

\subsection{El algoritmo UCB1}

El algoritmo \textbf{UCB1} (Auer et al., 2002) elige en cada tirada $t$ la
máquina $i$ que maximiza:

\begin{equation}
  \text{UCB1}(i) \;=\;
  \underbrace{\bar{x}_i}_{\text{explotación}}
  +
  \underbrace{\sqrt{\dfrac{2\ln n}{n_i}}}_{\text{exploración}}
\end{equation}

donde:
\begin{itemize}
  \item $\bar{x}_i = \frac{1}{n_i}\sum_{\text{tiradas en }i} r_t$: media
    empírica de recompensas de la máquina $i$.
  \item $n$: total de tiradas realizadas hasta el momento.
  \item $n_i$: número de veces que se ha tirado la máquina $i$.
\end{itemize}

\textbf{Interpretación geométrica del segundo término:}
El término $\sqrt{2\ln n / n_i}$ es proporcional al \emph{radio de un
intervalo de confianza} para $\mu_i$. Cuando $n_i$ es pequeño (máquina poco
explorada), este término es grande, incentivando visitar esa máquina. A medida
que $n_i$ crece, el intervalo se estrecha y confiamos más en la estimación
empírica.

\begin{tcolorbox}[teorema, title={Garantía de UCB1 (Auer et al., 2002)}]
Bajo las hipótesis H1 y H2, UCB1 garantiza un regret acumulado:
\[
  \mathbb{E}[R_n] \;\leq\; \sum_{i:\,\mu_i < \mu^*} \frac{8\ln n}{\Delta_i}
  + \left(1 + \frac{\pi^2}{3}\right)\sum_{i:\,\mu_i < \mu^*} \Delta_i
\]
donde $\Delta_i = \mu^* - \mu_i$ es la brecha de la máquina $i$ respecto a la
óptima. En particular, $\mathbb{E}[R_n] = O(\log n)$, que es óptimo.
\end{tcolorbox}

% ══════════════════════════════════════════════════════════════════════════════
\section{De bandidos a árboles de decisión}
% ══════════════════════════════════════════════════════════════════════════════

\subsection{El árbol de juego}

\hip{El problema puede modelarse como un \textbf{árbol de búsqueda} donde
cada nodo representa un estado del juego, cada arista representa una acción
legal, y existen estados terminales con recompensa definida ($+1$ victoria,
$0$ derrota, $0.5$ empate).}

\hip{El árbol es \textbf{demasiado grande} para explorarlo exhaustivamente. Por
esto no podemos aplicar minimax completo.}

\hip{El juego es de \textbf{información perfecta}: ambos jugadores conocen el
estado completo del sistema en todo momento. Esta hipótesis se relaja en
versiones avanzadas (ver POMCP).}

Formalmente, modelamos el problema como un \textbf{proceso de decisión de
Markov} (MDP):
\[
  \mathcal{M} = (\mathcal{S},\, \mathcal{A},\, T,\, R,\, \gamma)
\]
donde $\mathcal{S}$ es el espacio de estados, $\mathcal{A}$ el espacio de
acciones, $T: \mathcal{S}\times\mathcal{A} \to \Delta(\mathcal{S})$ la función
de transición, $R$ la función de recompensa, y $\gamma \in [0,1]$ el factor de
descuento.

\subsection{La conexión entre árboles y bandidos}

\begin{tcolorbox}[clave]
Cada nodo del árbol tiene varios hijos (movimientos posibles). Elegir cuál hijo
visitar es \textbf{exactamente el problema del bandido multi-brazo}: cada hijo
es una ``máquina'' cuya recompensa esperada es la calidad de ese movimiento,
que desconocemos.

La idea central de MCTS es \textbf{tratar cada nodo interno del árbol como una
instancia independiente del problema del bandido}.
\end{tcolorbox}

Esta observación, aparentemente simple, es la clave que hace funcionar MCTS.
Cada nodo almacena dos estadísticos:

\begin{itemize}
  \item $N(v)$: número de veces que el nodo $v$ ha sido visitado.
  \item $W(v)$: suma acumulada de recompensas obtenidas en simulaciones que
    pasaron por $v$.
\end{itemize}

El \textbf{valor estimado} del nodo es $Q(v) = W(v)/N(v)$.

% ══════════════════════════════════════════════════════════════════════════════
\section{El algoritmo MCTS: las cuatro fases}
% ══════════════════════════════════════════════════════════════════════════════

MCTS construye iterativamente un árbol $\mathcal{T}$ partiendo de la raíz
(estado actual del juego). Cada iteración ejecuta cuatro fases.

% ─── FASE 1 ──────────────────────────────────────────────────────────────────
\subsection{Fase 1 — Selección (UCT)}

\textbf{Hipótesis en uso: H1, H2, H3, H4.}

Partiendo de la raíz, recorremos el árbol hacia abajo usando el criterio
\textbf{UCT} (Upper Confidence Bounds for Trees, Kocsis \& Szepesvári, 2006),
que adapta UCB1 al contexto de árboles:

\begin{equation}
  \text{UCT}(v) \;=\;
  \underbrace{\dfrac{W(v)}{N(v)}}_{\text{explotación}}
  +\;
  c\underbrace{\sqrt{\dfrac{\ln N(\mathrm{padre}(v))}{N(v)}}}_{\text{exploración}}
  \label{eq:uct}
\end{equation}

donde $c > 0$ es el parámetro de exploración. El valor teóricamente motivado
es $c = \sqrt{2}$, aunque en la práctica se ajusta por validación cruzada.

\textbf{Procedimiento:} Desde la raíz, en cada nivel seleccionamos el hijo
con mayor $\text{UCT}(v)$, y descendemos hasta alcanzar un \textbf{nodo hoja}:
un nodo que tiene al menos un hijo no visitado, o un estado terminal.

\begin{tcolorbox}[nota]
\textbf{¿Por qué $\ln N(\mathrm{padre})$ y no $\ln n$ total?} Porque lo que
nos importa es el número de veces que hemos tenido la oportunidad de elegir
entre los hijos de ese nodo. Esto hace que el criterio sea localmente justo
en cada nivel del árbol.
\end{tcolorbox}

% ─── FASE 2 ──────────────────────────────────────────────────────────────────
\subsection{Fase 2 — Expansión}

\textbf{Hipótesis en uso: H3.}

Una vez en el nodo hoja, si el nodo no es terminal y tiene hijos no visitados,
añadimos \textbf{uno de esos hijos} al árbol, inicializado con $N = 0$,
$W = 0$.

\hip{Expandir \textbf{un único hijo} por iteración (en lugar de todos) es
suficiente para garantizar la convergencia y mantiene el coste por iteración
acotado por $O(b)$, donde $b$ es el factor de ramificación.}

En algunas implementaciones se expanden todos los hijos a la vez (expansión
total). Esto acelera la convergencia en árboles con branching factor pequeño,
pero puede ser prohibitivo si $b$ es grande.

% ─── FASE 3 ──────────────────────────────────────────────────────────────────
\subsection{Fase 3 — Simulación (Rollout)}

\textbf{Hipótesis en uso: H1, H2, H3, y las nuevas H7 y H8.}

Esta es la fase más característica y la que da el nombre «Monte Carlo» al
algoritmo.

\hip{Desde el nodo expandido, podemos \textbf{simular partidas} hasta el
final usando una política simple (típicamente movimientos uniformemente
aleatorios). El resultado de esa partida es una \textbf{estimación insesgada}
del valor real del nodo:
\[
  \mathbb{E}[\hat{G} \mid s] = V^{\pi_{\mathrm{ro}}}(s)
\]
donde $\pi_{\mathrm{ro}}$ es la política de rollout.}

\hip{La \textbf{varianza} de las recompensas es finita. Esto garantiza que el
promedio de múltiples rollouts converge al valor esperado por la ley de los
grandes números.}

\begin{tcolorbox}[nota]
La hipótesis H7 es \textbf{fuerte}: un rollout aleatorio es una estimación muy
ruidosa del valor real. Funciona porque se promedia sobre \emph{muchas}
iteraciones. En juegos donde la aleatoriedad no correlaciona con la victoria
(p.\,ej., Go a nivel profesional), el rollout aleatorio es casi inútil; esto
motivó reemplazarlo por redes neuronales en AlphaZero.
\end{tcolorbox}

El resultado es un escalar $\hat{G}$:
\begin{itemize}
  \item $\hat{G} = 1$: el jugador actual ganó la partida simulada.
  \item $\hat{G} = 0$: el jugador actual perdió.
  \item $\hat{G} = 0.5$: empate.
\end{itemize}

% ─── FASE 4 ──────────────────────────────────────────────────────────────────
\subsection{Fase 4 — Retropropagación (Backpropagation)}

\textbf{Hipótesis en uso: H3.}

El resultado $\hat{G}$ se propaga hacia arriba actualizando todos los nodos
del camino recorrido en la Fase 1:

\begin{equation}
  N(v) \;\leftarrow\; N(v) + 1, \qquad
  W(v) \;\leftarrow\; W(v) + \hat{G}
\end{equation}

para cada $v$ en el camino desde el nodo expandido hasta la raíz.

\begin{tcolorbox}[nota]
\textbf{Juegos de dos jugadores:} cuando $\hat{G}$ sube al nodo del oponente,
se usa $1 - \hat{G}$. Lo que es bueno para el jugador actual es malo para el
rival. Esta alternancia es automática si cada nodo almacena la perspectiva
del jugador que mueve en ese turno.
\end{tcolorbox}

% ─── DECISIÓN FINAL ──────────────────────────────────────────────────────────
\subsection{Decisión final}

Tras $n$ iteraciones, se elige la acción desde la raíz hacia el hijo con
\textbf{mayor $N(v)$} (más visitado). Alternativamente, se puede usar mayor
$Q(v) = W(v)/N(v)$.

\begin{tcolorbox}[clave]
¿Por qué más visitado en lugar de mayor $Q$? Porque el hijo más visitado es
aquel que UCT ha confirmado repetidamente como prometedor. Un hijo con $Q$
alto pero pocas visitas puede ser un \textbf{artefacto estadístico} (error de
muestreo con pocos datos).
\end{tcolorbox}

% ══════════════════════════════════════════════════════════════════════════════
\section{Garantías teóricas de convergencia}
% ══════════════════════════════════════════════════════════════════════════════

\textbf{Hipótesis en uso: H1, H2.}

\begin{tcolorbox}[teorema, title={Teorema de convergencia de UCT (Kocsis \& Szepesvári, 2006)}]
Sea $a^*$ la acción óptima en la raíz. Bajo las hipótesis H1 y H2 (recompensas
acotadas, i.i.d.\ por nodo), UCT satisface:
\[
  \Pr\!\left[\text{acción subóptima seleccionada tras } n
    \text{ iteraciones}\right] = O\!\left(\frac{\log n}{n}\right)
\]
En particular, la probabilidad de error tiende a $0$ cuando $n \to \infty$:
UCT \textbf{converge a la decisión óptima}.
\end{tcolorbox}

\textbf{Propiedades adicionales:}

\begin{itemize}
  \item \textbf{Anytime:} MCTS puede detenerse en cualquier momento y devuelve
    la mejor decisión encontrada hasta ese punto. Cuanto más tiempo se le da,
    mejor.
  \item \textbf{Sin función de evaluación:} no requiere conocimiento experto
    del dominio más allá de conocer las reglas del juego.
  \item \textbf{Asimetría:} MCTS explora el árbol de forma asimétrica,
    profundizando más en las ramas prometedoras y abandonando rápidamente las
    malas. A diferencia de minimax, no explora en anchura uniforme.
\end{itemize}

\textbf{Complejidad computacional por iteración:}

\begin{center}
\begin{tabular}{lc}
\toprule
\textbf{Fase} & \textbf{Complejidad} \\
\midrule
Selección (profundidad $d$) & $O(d)$ \\
Expansión (branching $b$) & $O(b)$ \\
Simulación (profundidad $h$) & $O(h)$ \\
Retropropagación & $O(d)$ \\
\midrule
\textbf{Total por iteración} & $O(d + b + h)$ \\
\textbf{Total ($n$ iteraciones)} & $O(n(d + b + h))$ \\
\bottomrule
\end{tabular}
\end{center}

% ══════════════════════════════════════════════════════════════════════════════
\section{Pseudocódigo comentado}
% ══════════════════════════════════════════════════════════════════════════════

\begin{lstlisting}[language=Python, caption={MCTS — pseudocódigo en Python}]
def MCTS(estado_raiz, n_iteraciones, c=1.41):
    """
    Ejecuta n_iteraciones del ciclo MCTS y retorna la mejor accion.
    c: parametro de exploracion (sqrt(2) por defecto).
    """
    raiz = Nodo(estado=estado_raiz, padre=None)

    for _ in range(n_iteraciones):

        # ── FASE 1: SELECCION ─────────────────────────────────────────
        # Hipotesis: usamos UCT para navegar el arbol ya construido.
        # Descendemos hasta un nodo hoja (no completamente expandido
        # o estado terminal).
        nodo = raiz
        while nodo.completamente_expandido() and not nodo.es_terminal():
            nodo = max(nodo.hijos,
                       key=lambda v: v.W/v.N + c*(log(v.padre.N)/v.N)**0.5)

        # ── FASE 2: EXPANSION ─────────────────────────────────────────
        # Hipotesis: hay hijos no visitados; expandimos uno.
        if not nodo.es_terminal():
            nodo = nodo.expandir_hijo_no_visitado()

        # ── FASE 3: SIMULACION (ROLLOUT) ──────────────────────────────
        # Hipotesis: el rollout aleatorio estima el valor real del nodo
        # (estimacion insesgada pero ruidosa).
        resultado = simular_partida_aleatoria(nodo.estado)

        # ── FASE 4: RETROPROPAGACION ──────────────────────────────────
        # Actualizamos N y W en todos los nodos del camino recorrido.
        while nodo is not None:
            nodo.N += 1
            nodo.W += resultado
            resultado = 1 - resultado  # alternamos perspectiva (2 jugadores)
            nodo = nodo.padre

    # Elegir la mejor accion: hijo de la raiz con mayor N (mas visitado)
    return max(raiz.hijos, key=lambda v: v.N).accion
\end{lstlisting}

\textbf{Estructura del nodo:}
\begin{lstlisting}[language=Python, caption={Clase Nodo para MCTS}]
class Nodo:
    def __init__(self, estado, padre, accion=None):
        self.estado  = estado       # estado del juego en este nodo
        self.padre   = padre        # nodo padre (None si es raiz)
        self.accion  = accion       # accion que llevo aqui desde el padre
        self.hijos   = []           # lista de nodos hijo expandidos
        self.N = 0                  # numero de visitas
        self.W = 0.0                # suma acumulada de recompensas

    def completamente_expandido(self):
        acciones_legales = self.estado.acciones_legales()
        acciones_expandidas = [h.accion for h in self.hijos]
        return set(acciones_legales) == set(acciones_expandidas)

    def expandir_hijo_no_visitado(self):
        acciones_no_visitadas = [
            a for a in self.estado.acciones_legales()
            if a not in [h.accion for h in self.hijos]
        ]
        accion = random.choice(acciones_no_visitadas)
        nuevo_estado = self.estado.aplicar(accion)
        hijo = Nodo(estado=nuevo_estado, padre=self, accion=accion)
        self.hijos.append(hijo)
        return hijo
\end{lstlisting}

% ══════════════════════════════════════════════════════════════════════════════
\section{Cuándo se cumplen y cuándo se violan las hipótesis}
% ══════════════════════════════════════════════════════════════════════════════

\begin{longtable}{L{3cm} L{5cm} L{6cm}}
\toprule
\textbf{Hipótesis} & \textbf{¿Cuándo se viola?} & \textbf{Consecuencia y solución} \\
\midrule
\endhead
\textbf{H1} Recompensas i.i.d. &
  Oponente adaptativo; entornos no estacionarios; juegos con elementos
  aleatorios dependientes del tiempo &
  UCT puede no converger. Solución: usar variantes robustas o modelar
  explícitamente la estocasticidad. \\[6pt]

\textbf{H2} Recompensas acotadas &
  Recompensas continuas sin límite (p.\,ej., puntuación acumulable
  indefinidamente) &
  Las cotas de regret no aplican directamente. Solución: normalizar las
  recompensas al rango $[0,1]$. \\[6pt]

\textbf{H3} Árbol de juego &
  Juegos con ciclos (tablero infinito, posiciones repetidas en ajedrez/Go) &
  El árbol se convierte en un grafo (DAG). Solución: tablas de transposición
  (tabla hash de estados $\to$ estadísticos). \\[6pt]

\textbf{H5} Información perfecta &
  Póker, bridge, Starcraft, robótica con sensores ruidosos &
  El estado real es latente (POMDP). Solución: POMCP (opera sobre espacios
  de creencias) o IS-MCTS (Information Set MCTS). \\[6pt]

\textbf{H6} Expansión de un hijo &
  Branching factor muy pequeño ($b \leq 5$) &
  La expansión total (todos los hijos a la vez) puede ser más eficiente. \\[6pt]

\textbf{H7} Rollout estima bien &
  Juegos donde el azar aleatorio no correlaciona con la victoria
  (Go a nivel profesional, ajedrez) &
  El rollout es inútil o contraproducente. Solución: reemplazarlo por una
  red neuronal de valor (AlphaZero). \\[6pt]

\textbf{H8} Varianza finita &
  Recompensas de cola pesada (distribuciones de Pareto) &
  La ley de los grandes números converge muy lentamente. Solución:
  algoritmos de bandido robustos (mediana de medias). \\

\bottomrule
\caption{Resumen de hipótesis, condiciones de violación y soluciones.}
\end{longtable}

% ══════════════════════════════════════════════════════════════════════════════
\section{Recursos por nivel de profundidad}
% ══════════════════════════════════════════════════════════════════════════════

\subsection*{Nivel 0 — Intuición visual (sin matemáticas)}

\begin{itemize}
  \item \href{https://medium.com/data-science/the-animated-monte-carlo-tree-search-mcts-c05bb48b018c}%
    {\textbf{The Animated MCTS} — Thomas Kurbiel, Medium}: animaciones GIF de
    cada una de las cuatro fases. Ideal para construir la intuición visual antes
    de abordar la fórmula.

  \item \href{https://medium.com/@_michelangelo_/monte-carlo-tree-search-mcts-algorithm-for-dummies-74b2bae53bfa}%
    {\textbf{MCTS for Dummies} — Medium}: analogía con exploración de mapas y
    lenguaje completamente accesible.
\end{itemize}

\subsection*{Nivel 1 — Introductorio con código Python}

\begin{itemize}
  \item \href{https://int8.io/monte-carlo-tree-search-beginners-guide/}%
    {\textbf{Int8 — Beginner's Guide to MCTS}}: el mejor recurso introductorio
    completo. Cubre teoría, pseudocódigo e implementación en Python con juego
    de ejemplo. \textbf{Altamente recomendado como primer recurso técnico}.

  \item \href{https://ai-boson.github.io/mcts/}%
    {\textbf{ai-boson — MCTS Python Step by Step}}: implementación limpia y
    comentada con Tic-Tac-Toe. Ideal para leer en paralelo al código.
\end{itemize}

\subsection*{Nivel 2 — Introductorio con teoría básica}

\begin{itemize}
  \item \href{https://towardsdatascience.com/monte-carlo-tree-search-an-introduction-503d8c04e168/}%
    {\textbf{TDS — MCTS An Introduction} — Benjamin Wang}: narrativa clara con
    fórmulas básicas y ejemplos de aplicación.

  \item \href{https://gibberblot.github.io/rl-notes/single-agent/mcts.html}%
    {\textbf{Mastering Reinforcement Learning — MCTS Chapter}}: libro de texto
    libre en línea con MCTS aplicado a RL, incluyendo GridWorld y Freeway.
    Tiene visualizaciones interactivas.

  \item \href{https://lilianweng.github.io/posts/2018-01-23-multi-armed-bandit/}%
    {\textbf{Lilian Weng — The Multi-Armed Bandit Problem}}: el mejor post para
    el prerrequisito de bandidos. Con gráficas, código y pruebas informales.
\end{itemize}

\subsection*{Nivel 3 — Tutorial académico completo}

\begin{itemize}
  \item \href{https://www.informs-sim.org/wsc18papers/includes/files/021.pdf}%
    {\textbf{WSC 2018 — Monte Carlo Tree Search: A Tutorial (PDF)}}: artículo
    tutorial publicado en el Winter Simulation Conference. Riguroso pero
    accesible. Cubre historia, variantes y aplicaciones.

  \item \href{https://courses.cs.washington.edu/courses/cse599i/18wi/resources/lecture19/lecture19.pdf}%
    {\textbf{UWashington — Lecture 19: MCTS (PDF)}}: diapositivas de clase
    universitaria con demostraciones informales del teorema de convergencia.

  \item \href{https://uq.pressbooks.pub/mastering-reinforcement-learning/chapter/monte-carlo-tree-search/}%
    {\textbf{UQ Press — MCTS en RL}}: versión en libro de texto con contexto
    de aprendizaje por refuerzo completo.
\end{itemize}

\subsection*{Nivel 4 — Artículos de investigación de referencia}

\begin{itemize}
  \item \href{https://link.springer.com/article/10.1007/s10462-022-10228-y}%
    {\textbf{Springer (2022) — MCTS: A Review of Recent Modifications}}: revisión
    sistemática de todas las variantes modernas de MCTS (PUCT, RAVE, GRAVE,
    progressive widening, etc.) con tabla comparativa de aplicaciones.

  \item \href{https://dke.maastrichtuniversity.nl/m.winands/documents/Encyclopedia_MCTS.pdf}%
    {\textbf{Winands — Encyclopedia of MCTS (PDF)}}: entrada de enciclopedia
    académica. La referencia más exhaustiva del estado del arte clásico.

  \item \href{https://www.chessprogramming.org/UCT}%
    {\textbf{Chessprogramming Wiki — UCT}}: detalles de implementación de UCT
    específicos para juegos de tablero, con variantes y trucos de rendimiento.
\end{itemize}

\subsection*{Prerrequisito: problema del bandido multi-brazo}

\begin{itemize}
  \item \href{https://lilianweng.github.io/posts/2018-01-23-multi-armed-bandit/}%
    {Lilian Weng — Multi-Armed Bandit (ya citado, es el mejor)}

  \item \href{https://www.turing.com/kb/guide-on-upper-confidence-bound-algorithm-in-reinforced-learning}%
    {\textbf{Turing.com — Guía de UCB}}: explicación paso a paso del algoritmo
    UCB1 con ejemplos numéricos.

  \item \href{https://cse442-17f.github.io/LinUCB/}%
    {\textbf{CSE 442 — Multi-Armed Bandits (interactivo)}}: visualización
    interactiva del dilema exploración-explotación en bandidos.
\end{itemize}

% ══════════════════════════════════════════════════════════════════════════════
\section{Resumen: el mapa conceptual del algoritmo}
% ══════════════════════════════════════════════════════════════════════════════

\begin{center}
\begin{tabular}{ccccc}
  \toprule
  \textbf{Fase} & \textbf{Pregunta} & \textbf{Herramienta} & \textbf{Hipótesis clave} & \textbf{Parámetro} \\
  \midrule
  Selección & ¿Por qué rama bajar? & UCT & H1, H2 & $c$ \\
  Expansión & ¿Qué nodo añadir? & Aleatorio o heurístico & H6 & $n_{\min}$ \\
  Simulación & ¿Cuál es el valor? & Rollout aleatorio & H7, H8 & política $\pi_{\mathrm{ro}}$ \\
  Retropropagación & ¿Cómo actualizar? & Promedio acumulado & H3 & — \\
  \bottomrule
\end{tabular}
\end{center}

\vspace{1em}

\begin{tcolorbox}[clave, title={Resumen en una frase}]
MCTS trata cada nodo del árbol de juego como un problema de bandido
multi-brazo, usa UCT para balancear exploración y explotación en cada nivel,
y estima el valor de los estados desconocidos mediante simulaciones
Monte Carlo — sin necesitar ninguna función de evaluación escrita a mano.
\end{tcolorbox}

% ══════════════════════════════════════════════════════════════════════════════
\section*{Referencias principales}
% ══════════════════════════════════════════════════════════════════════════════
\addcontentsline{toc}{section}{Referencias principales}

\begin{enumerate}
  \item Kocsis, L., \& Szepesvári, C. (2006). \textit{Bandit based Monte-Carlo
    Planning}. In \textit{European Conference on Machine Learning (ECML)},
    282--293. \textbf{[Artículo fundacional de UCT]}

  \item Auer, P., Cesa-Bianchi, N., \& Fischer, P. (2002). Finite-time analysis
    of the multiarmed bandit problem. \textit{Machine Learning}, 47, 235--256.
    \textbf{[Garantías de UCB1]}

  \item Lai, T.\,L., \& Robbins, H. (1985). Asymptotically efficient adaptive
    allocation rules. \textit{Advances in Applied Mathematics}, 6(1), 4--22.
    \textbf{[Cota inferior del regret]}

  \item Browne, C. et al. (2012). A survey of Monte Carlo tree search methods.
    \textit{IEEE TCIAIG}, 4(1), 1--43. \textbf{[La survey de referencia clásica]}

  \item Silver, D. et al. (2016). Mastering the game of Go with deep neural
    networks and tree search. \textit{Nature}, 529, 484--489.
    \textbf{[AlphaGo]}

  \item Silver, D. et al. (2018). A general reinforcement learning algorithm
    that masters chess, shogi, and Go through self-play. \textit{Science},
    362(6419), 1140--1144. \textbf{[AlphaZero]}
\end{enumerate}

\end{document}
